% Independent addition by /root/frobenius, after C.4h--l.
% New keys are in C4_prym_references.bib; other keys are already in C4 references.
\subsection{A Prym improvement: only a possible factor of two for surfaces}
\label{sec:c4-prym-improvement}

\paragraph{C.4m (Proved: twice the identity for principally polarized varieties of dimension at most five).}
Let $(B,\Xi)$ be a principally polarized abelian variety of dimension
$1\le h\le5$ over an algebraically closed field of arbitrary characteristic.
Then
\begin{equation}\label{eq:c4-ppav-five-two}
                         2\in I_J(B).
\end{equation}
More precisely, there is an integral curve cycle $\alpha$ on $B$ with
\begin{equation}\label{eq:c4-twice-minimal-cycle}
       \operatorname{cl}(\alpha)
          =2\,\frac{\xi^{h-1}}{(h-1)!},\qquad \xi=c_1(\Xi),
\end{equation}
in rational $\ell$-adic cohomology for every $\ell$ different from the
characteristic. We only need this equality for one such $\ell$ to obtain
the equality of homomorphisms in \eqref{eq:c4-ppav-five-two}.

\emph{Proof in characteristic different from two.}
For $h\le3$ we can use the product-of-Jacobians theorem
\citep[Theorem 4, p.~2 of the opened preprint]{cyOortUeno1973}.
For $h=4,5$, \citet[Theorem 6.5(i), p.~178]{frC4Beauville1977}
realizes $(B,\Xi)$ as a generalized Prym. Let $\widetilde C$ be its
nodal covering curve, with smooth normalization components
$C_1,\ldots,C_t$. The composition
\[
 j:B\longrightarrow J(\widetilde C)
            \longrightarrow J:=\prod_{a=1}^{t}J(C_a)
\]
is a homomorphism of abelian varieties. The generalized Jacobian in
the middle is allowed to be semiabelian; the target $J$ is abelian.
Proposition 2.2, p.~155, and Theorem 5.4, pp.~175--176, of the same
paper give
\[
                         j^*\Theta_J\equiv_{\mathrm{alg}}2\Xi.
\]
Indeed, the first statement identifies the generalized theta class
with pullback of the product theta class of the normalization, and
the second identifies its restriction to the Prym.
Writing $\lambda_B$ and $\lambda_J$ for the principal polarizations,
we define the \emph{integral} polarized adjoint
\[
             q=j^\dagger:=\lambda_B^{-1}\widehat j\lambda_J:J\to B.
\]
The pullback identity gives
\begin{equation}\label{eq:c4-prym-adjoint}
 qj=\lambda_B^{-1}(\widehat j\lambda_Jj)
    =\lambda_B^{-1}(2\lambda_B)=[2]_B.
\end{equation}
We have inverted only principal polarizations, which are actual
isomorphisms. In particular, a possible finite kernel of $j$ has not
been divided out.

We also record the curve needed for specialization. Choose Abel maps
$\iota_a:C_a\to J(C_a)$, and let $q_a$ be the restriction of $q$ to
the $a$th factor. Put
\[
                 \alpha=\sum_a(q_a\iota_a)_*[C_a].
\]
Components mapped to a point contribute zero. The symmetric map
$\widehat B\to B$ associated with this curve class is
\[
 \sum_aq_a\lambda_a^{-1}\widehat{q_a}
   =q\lambda_J^{-1}\widehat q
   =\lambda_B^{-1}\widehat j\lambda_Jj\lambda_B^{-1}
   =2\lambda_B^{-1}.
\]
Under Poincar\'e duality, the class of a curve is the alternating
pairing obtained by integrating products of two degree-one classes
over that curve. The preceding equality says that this pairing is
twice the inverse principal-polarization pairing. In a symplectic
basis, the complementary exterior power
$\xi^{h-1}/(h-1)!$ gives precisely that inverse pairing.
Since degree-one products span $H^2(B,\mathbb Q_\ell)$, the pairing
determines the curve class. This proves
\eqref{eq:c4-twice-minimal-cycle}.

\emph{Characteristic two.}
By \citet[Theorem 2.4.1, p.~286, and Corollary 2.4.2 with its proof,
p.~287]{frC4Oort1971}, a principally polarized abelian variety lifts,
with its polarization, to an abelian scheme $\mathcal B$ over
$W(k)$. A principal polarization is separable even in characteristic
two, so the stated hypothesis is satisfied. On its geometric generic
fibre the preceding construction supplies an integral curve cycle
with class $2\xi^{h-1}/(h-1)!$.
We use only the separability-implies-smoothness direction of
Theorem 2.4.1: the converse in its printed statement is explicitly
retracted in \citet[p.~3 of the opened preprint]{cyOortUeno1973}.

The cycle and its finitely many components descend to a finite
extension of the fraction field. After the corresponding finite
extension of complete discrete valuation rings, we take the closures
of those components and their special-fibre cycles. This gives an
integral cycle $\alpha_0$ on $B$; the residue field remains $k$ because
it is algebraically closed. The specialization map commutes with
the $\ell$-adic cycle class map for $\ell\ne2$
\citep[Proposition 20.3, p.~399, and Example 20.3.5,
p.~400]{frC4Fulton1998}. Hence
\[
                 \operatorname{cl}(\alpha_0)
                    =2\,\frac{\xi^{h-1}}{(h-1)!}.
\]
In particular, the same equality holds numerically. Applying the
integral Matsusaka construction in
\eqref{eq:c4-matsusaka-integer}, with $n=2$ and $\chi(\Xi)=1$,
factors $[2]_B$ through the Jacobians of the normalizations of the
components of $\alpha_0$. This is an equality of algebraic
homomorphisms in characteristic two. We have specialized an actual
integral cycle; no inference from a crystalline realization, and
no inversion of two, enters the argument.\hfill$\square$

\paragraph{C.4n (Proved: the surface obstruction divides two, in every characteristic).}
For any abelian surface $A$ over an algebraically closed field,
\begin{equation}\label{eq:c4-surface-two}
 \boxed{\qquad s_J(A)\mid\gcd(2,\delta(A)).\qquad}
\end{equation}
Consequently every surface admitting an ample line bundle of odd
Euler characteristic is an actual direct summand of a product of
Jacobians. In particular, for a complex surface whose
N\'eron--Severi group is generated by a polarization of type $(1,d)$,
every odd $d$ satisfies Question 2 affirmatively.

\emph{Proof.}
We first make a useful consequence of C.4m explicit. If $X$ has
dimension four or five, principalizing any ample $L$ gives an
isogeny $f:X\to B$ of degree $\chi(L)$ to a principally polarized
variety of the same dimension. This is the all-characteristic
group-scheme construction used in C.4j
\citep[Section 23, Theorem 4 and Corollary 1,
pp.~233--234]{frC4Mumford1970}.
By \eqref{eq:c4-isogeny-divisibility} and C.4m,
$2\chi(L)\in I_J(X)$. Taking the gcd over ample $L$, and using
the equality of the ample and unrestricted gcds proved in C.4j,
we obtain
\begin{equation}\label{eq:c4-four-five-delta}
              s_J(X)\mid2\delta(X)\qquad(\dim X=4,5).
\end{equation}

Now take $X=A\times\widehat A$. The Poincar\'e line bundle is
unimodular and has Euler characteristic $1$ on this fourfold.
Therefore $\delta(X)=1$. Since $A$ is an actual direct factor of
$X$, restriction and projection of a factorization on $X$ show
$I_J(X)\subset I_J(A)$. Equation \eqref{eq:c4-four-five-delta}
therefore gives $2\in I_J(A)$.
C.4j also gives $\delta(A)\in I_J(A)$, proving
\eqref{eq:c4-surface-two}. If $\delta(A)$ is odd, an integral
B\'ezout combination of the two factorizations gives $1\in I_J(A)$;
the explicit splitting maps in C.4h then give the required actual
direct summand. For a Picard-rank-one surface of type $(1,d)$,
C.4k gives $\delta(A)=d$.\hfill$\square$

For a finite construction of the factorization on $X$, choose an
ample $H$, put $p(t)=\chi(\mathcal P_A\otimes H^t)$, and take
$T\ge1$ so large that the five line bundles with
$t=T,T+1,\ldots,T+4$ are ample. The integral interpolation identity
\[
 1=p(0)=\sum_{j=0}^{4}(-1)^j
       \binom{T+j-1}{j}\binom{T+4}{4-j}\,p(T+j)
\]
expresses $2$ as an integral combination of five available
factorizations of $2p(T+j)$. The identity follows by evaluating the
degree-at-most-four Lagrange interpolation formula at zero; its
displayed coefficients are integers. Thus the gcd step does not
require an infinite construction.

There is also a construction using just one principalization of
$A\times\widehat A$, besides the factorial factorization.
Fix an ample $H$ on $X$, and take $N$ sufficiently large and divisible
by $3\cdot4!=72$. Then
$L=\mathcal P_A\otimes H^N$ is ample and
\[
 \chi(L)
 =1+\sum_{r=1}^{4}\frac{N^r}{r!(4-r)!}
             \deg\bigl(c_1(\mathcal P_A)^{4-r}c_1(H)^r\bigr)
 \equiv1\pmod3.
\]
Every displayed coefficient after the constant is an integer
divisible by $3$. If $m=\chi(L)$, principalization and C.4m give
$2m\in I_J(A)$, whereas C.4i gives $6\in I_J(A)$.
The equality $\gcd(6,2m)=2$ provides the factorization of $[2]$.
This version makes the removal of the possible factor of three
particularly concrete.
